\subsubsection{Subproceso 1: Acreditaci\'on y Pre-requisitos}

El subproceso se activa d\'ias antes de la semana de matr\'icula con dos actividades en paralelo. En el frente acad\'emico, los Departamentos y Facultades cierran las actas de evaluaci\'on del ciclo previo sin dejar reclamos pendientes (RR-0570-2022, Art.~6\textdegree{}). Las Escuelas Profesionales, a su vez, remiten a la DIRCE las n\'ominas de estudiantes pr\'oximos a egresar para su turno preferencial (Art.~42\textdegree{}). En el frente administrativo, el estudiante regular genera su orden de pago en el portal Intralu y cancela el aporte al Fondo de Autoseguro Estudiantil (RR-2461-2023, nums.~6.1.1 y 6.2)~\cite{autoseguro2023}. Adem\'as subsana posibles deudas de libros o materiales en su Facultad~\cite{tupa2024}. Si el estudiante registra cursos desaprobados por segunda o tercera vez, debe recibir previamente la consejer\'ia de su docente tutor. Este luego autoriza la activaci\'on de su c\'odigo en la plataforma (Arts.~8\textdegree{} y 48\textdegree{}).

Con las actas cerradas y los expedientes recibidos, el sistema SIGA-ORCE realiza el cruce autom\'atico de recaudaci\'on y verifica los bloqueos reportados. Cuando el pago del autoseguro est\'a confirmado, el sistema calcula el promedio ponderado de los dos \'ultimos semestres regulares del alumno. Define con ello su tope m\'aximo de cr\'editos y establece su puesto en el orden de m\'erito (Art.~9\textdegree{}, inc.~a). A partir de este orden, la DIRCE y la Comisi\'on de Matr\'icula estructuran los turnos de atenci\'on de la Facultad. Los tres bloques son: estudiantes en riesgo acad\'emico, estudiantes de promoci\'on y grupos regulares ordenados de mayor a menor promedio.

Como resultado de este procesamiento, el sistema genera la Ficha de Matr\'icula personalizada con el plan de estudios, especialidad y l\'imite de cr\'editos de cada alumno. Luego, la Intranet DIRCE publica el cronograma general de turnos de la Facultad y muestra al estudiante los cuatro controles de habilitaci\'on (\textit{Grupo de matr\'icula}, \textit{Bloqueos}, \textit{Adeudos} y \textit{Pagos}). Al encontrarse estos requisitos validados, el estudiante queda formalmente acreditado y habilitado. Con esto concluye el Subproceso~1 y el estudiante espera la apertura de su turno para iniciar la selecci\'on de asignaturas en el Subproceso~2.

El an\'alisis del proceso actual en el periodo 2026-2 evidencia importantes deficiencias. Son operativas y tecnol\'ogicas:
\begin{enumerate}
    \item \textbf{Inoperatividad del filtro de deudas:} El TUPA (PA5400954D)~\cite{tupa2024} y el Reglamento de Matr\'icula (Arts.~46\textdegree{} y 47\textdegree{})~\cite{reglamento2022} exigen no registrar deudas con la Facultad. En la pr\'actica, el sistema no contrast\'o los adeudos de bibliotecas o laboratorios. La habilitaci\'on dependi\'o exclusivamente de haber pagado el autoseguro estudiantil.
    \item \textbf{Falta de disponibilidad y ca\'idas de los portales:} Se constat\'o la interrupci\'on nocturna del servicio en los portales Intralu e Intranet DIRCE a partir de las 11:00~p.\,m. (\texttt{ERR\_CONNECTION\_REFUSED}). Esta medici\'on se obtuvo en periodo regular, fuera de la semana de matr\'icula. Aun as\'i, es recurrente que durante las jornadas de inscripci\'on la plataforma sufra ca\'idas o demoras severas por sobrecarga concurrente de usuarios.
    \item \textbf{Carga manual y as\'incrona de tr\'amites previos:} Las fichas se arman con la base de datos congelada del ciclo anterior. Si un alumno tiene convalidaciones, exoneraciones o un cambio de plan en curso, la informaci\'on debe viajar f\'isicamente desde la Facultad hasta la DIRCE para su digitaci\'on manual. Esta demora impide que el estudiante tenga su ficha actualizada al momento de abrirse su turno de matr\'icula.
    \item \textbf{Fallas en la configuraci\'on de promociones y planes:} Durante el levantamiento se identificaron errores en casos reales. El primero fue el cambio de plan a estudiantes de Ingenier\'ia F\'isica, que dej\'o cursos aprobados sin convalidar en la ficha. El segundo fue la omisi\'on del turno preferencial a alumnos egresantes de la escuela de F\'isica, a quienes el sistema envi\'o indebidamente a turnos regulares.
\end{enumerate}

Esta descripci\'on del proceso actual sirve de base para el diagn\'ostico de causas ra\'iz y la propuesta de mejora tecnol\'ogica que se presentan en los siguientes cap\'itulos.
